% =============================================================================
%                    WEEK 1: INTRODUCTION, NAVIGATION & GEOMETRY
% =============================================================================
\section{Week 1: Introduction to Autonomous Robots, Simple Navigation \& Geometry Representation}

\subsection{Topic Summary}

\begin{keybox}[Learning Objectives]
By the end of this week, you should be able to:
\begin{itemize}
  \item Explain the core challenges in building autonomous robots (sense--plan--act cycle).
  \item Describe and analyse Bug algorithms (Bug 0, Bug 1, Bug 2, Tangent Bug).
  \item Prove completeness and derive path length bounds for Bug algorithms.
  \item Represent workspace geometry using polygons, primitives, and point clouds.
  \item Distinguish between convex and non-convex sets and their representations.
\end{itemize}
\end{keybox}

\textbf{Big Picture.} This week provides the foundations for robot autonomy. We start with what makes robotics hard (uncertainty, limited sensing, computational cost), then introduce Bug algorithms---simple reactive navigation strategies with provable guarantees. Finally, we cover how to mathematically represent the robot's environment.

\textbf{What Examiners Like to Ask:}
\begin{itemize}
  \item ``Trace Bug 1/Bug 2 on this environment and identify hit/leave points.''
  \item ``Derive upper/lower bounds on path length for Bug 1 or Bug 2.''
  \item ``Prove that Bug 1 is complete.''
  \item ``Compare Bug 1 vs Bug 2---when does each perform better?''
  \item ``Represent this shape using half-planes (primitives).''
  \item ``Define convexity and give an example of a convex hull.''
\end{itemize}

\bigseparator

\subsection{Core Concepts}

\textbf{Roadmap:}
\begin{enumerate}
  \item What is a robot? The sense--plan--act cycle
  \item Bug algorithms for point robot navigation
  \item Geometry representation methods
\end{enumerate}

% -----------------------------------------------------------------------------
\subsubsection{What is an Autonomous Robot?}

\begin{defbox}[Robot]
A \term{robot} is a physical system that \textbf{observes} the state of the world and \textbf{changes} the state of the world.
\end{defbox}

\begin{defbox}[Autonomous Robot]
An \term{autonomous robot} performs tasks without continuous human intervention, using its own sensing, planning, and actuation.
\end{defbox}

\textbf{The Sense--Plan--Act Cycle:}
\begin{center}
\begin{tabular}{@{} l p{10cm} @{}}
\toprule
\textbf{Component} & \textbf{Description} \\
\midrule
Sense & Perceive the world using sensors (cameras, LIDAR, tactile sensors). \\
Plan & Decide what actions to take to achieve a goal. \\
Act & Execute actions using actuators (motors, grippers). \\
\bottomrule
\end{tabular}
\end{center}

\textbf{Key Challenges in Robotics:}
\begin{itemize}
  \item \textbf{Limited knowledge:} Sensors are noisy, incomplete, and have finite range.
  \item \textbf{Computational cost:} Planning can be exponentially hard.
  \item \textbf{Hardware limitations:} Control is imperfect; actuators have uncertainty.
\end{itemize}

\separator

% -----------------------------------------------------------------------------
\subsubsection{Motion Planning: The Piano Mover's Problem}

\begin{defbox}[Motion Planning Problem]
Given a robot, a start configuration, a goal configuration, and a description of obstacles, find a \textbf{collision-free path} from start to goal (or report that none exists).
\end{defbox}

\begin{tipbox}[Exam Tip: Complexity]
Motion planning is computationally hard in general:
\begin{itemize}
  \item Sofa Mover (3 DOF): $O(n^{2+\epsilon})$
  \item Piano Mover (6 DOF): Polynomial but no practical algorithm known
  \item $n$ Disks in Plane: NP-hard
  \item Generalized Mover: PSPACE-complete
\end{itemize}
This motivates the use of heuristic and sampling-based algorithms.
\end{tipbox}

\bigseparator

% =============================================================================
%                              BUG ALGORITHMS
% =============================================================================
\subsection{Bug Algorithms}

Bug algorithms are simple, reactive navigation strategies for a \textbf{point robot} with:
\begin{itemize}
  \item Known direction to goal (can measure distance $d(x, y)$ between any two points).
  \item Local sensing only (tactile/contact sensor or finite-range sensor).
  \item No global map of obstacles.
\end{itemize}

\begin{defbox}[Workspace Notation]
\begin{itemize}
  \item $\workspace \subseteq \mathbb{R}^2$: the \term{workspace} (environment).
  \item $\mathcal{WO}_i$: obstacle $i$ in the workspace.
  \item $\workspace_{\text{free}} = \workspace \setminus \bigcup_i \mathcal{WO}_i$: the \term{free workspace}.
\end{itemize}
\end{defbox}

\textbf{Assumptions for Bug algorithms:}
\begin{enumerate}
  \item Perfect localization (robot knows its position at all times).
  \item Contact sensor available.
  \item Robot can measure distance to goal.
  \item Finitely many obstacles; a line intersects each obstacle finitely many times.
  \item Workspace is bounded: $\workspace \subset B_r(x)$ for some finite $r$.
\end{enumerate}

\separator

% --- Bug 0 ---
\subsubsection{Bug 0 Algorithm}

\begin{thmbox}[Bug 0 Strategy]
\begin{enumerate}
  \item Head directly toward the goal.
  \item If an obstacle is encountered, follow its boundary (fixed direction: left or right) until you can head toward the goal again.
  \item Repeat.
\end{enumerate}
\end{thmbox}

\begin{tipbox}[Bug 0 Failure Case]
Bug 0 is \textbf{not complete}. It can get trapped in certain environments.

\textbf{Example:} A left-turning robot facing a ``spiral'' obstacle that always redirects it away from the goal.
\end{tipbox}

\separator

% --- Bug 1 ---
\subsubsection{Bug 1 Algorithm}

Bug 1 adds \textbf{memory} to avoid the traps that defeat Bug 0.

\begin{thmbox}[Bug 1 Strategy]
\begin{enumerate}
  \item From current position, head toward $q_{\text{goal}}$.
  \item If obstacle encountered at \term{hit point} $q_H^i$:
  \begin{enumerate}
    \item \textbf{Circumnavigate} the entire obstacle boundary.
    \item Record the point $q_L^i$ (\term{leave point}) closest to the goal.
    \item Return to $q_L^i$ by wall-following.
  \end{enumerate}
  \item From $q_L^i$, head toward goal. If this moves into the obstacle, report \textbf{failure}.
  \item Repeat until goal reached or failure.
\end{enumerate}
\end{thmbox}

\begin{defbox}[Hit Point and Leave Point]
\begin{itemize}
  \item $q_H^i$: the \term{hit point}---where the robot first contacts obstacle $i$.
  \item $q_L^i$: the \term{leave point}---the point on obstacle $i$'s boundary closest to the goal.
\end{itemize}
\end{defbox}

\textbf{Bug 1 Path Length Bounds:}

Let $D$ = straight-line distance from $q_{\text{start}}$ to $q_{\text{goal}}$, and $P_i$ = perimeter of obstacle $i$.

\begin{center}
\renewcommand{\arraystretch}{1.3}
\begin{tabular}{@{} l l @{}}
\toprule
\textbf{Bound} & \textbf{Expression} \\
\midrule
Lower bound & $D$ \\
Upper bound & $D + \dfrac{3}{2} \sum_i P_i$ \\
\bottomrule
\end{tabular}
\end{center}

\textbf{Why $\frac{3}{2}$?} In the worst case, the robot circumnavigates the full obstacle ($P_i$), then backtracks up to half the perimeter to reach the leave point ($\frac{1}{2}P_i$).

\begin{thmbox}[Bug 1 Completeness]
\textbf{Theorem:} Bug 1 is \term{complete}---it will find a path in finite time if one exists, or correctly report failure.
\end{thmbox}

\begin{exbox}[Proof Sketch: Bug 1 Completeness]
Suppose Bug 1 is incomplete. Then either:
\begin{enumerate}
  \item \textbf{Never terminates:} But each leave point is strictly closer to the goal than its corresponding hit point. Since there are finitely many obstacles, there are finitely many hit/leave pairs. Eventually, the robot reaches the goal.
  \item \textbf{Terminates incorrectly:} The leave point $q_L^i$ is the closest point on the boundary to the goal. If moving toward the goal from $q_L^i$ re-enters the obstacle, then by the \textbf{Jordan Curve Theorem}, the line to the goal must cross the boundary an even number of times, implying another boundary point is closer---contradicting the definition of $q_L^i$.
\end{enumerate}
\end{exbox}

\separator

% --- Bug 2 ---
\subsubsection{Bug 2 Algorithm}

Bug 2 is a \textbf{greedy} algorithm that uses the \term{m-line} (the line from $q_{\text{start}}$ to $q_{\text{goal}}$).

\begin{thmbox}[Bug 2 Strategy]
\begin{enumerate}
  \item Head toward goal along the \term{m-line}.
  \item If obstacle encountered at hit point $q_H^i$:
  \begin{enumerate}
    \item Follow the obstacle boundary.
    \item Leave when you re-encounter the m-line at a point \textbf{closer to the goal} than $q_H^i$ (and the path to goal is unimpeded).
  \end{enumerate}
  \item If you return to $q_H^i$ without finding such a point, report \textbf{failure}.
  \item Repeat until goal reached.
\end{enumerate}
\end{thmbox}

\begin{defbox}[M-Line]
The \term{m-line} is the straight line segment connecting $q_{\text{start}}$ to $q_{\text{goal}}$.
\end{defbox}

\textbf{Bug 2 Path Length Bounds:}

Let $n_i$ = number of times the m-line intersects obstacle $i$.

\begin{center}
\renewcommand{\arraystretch}{1.3}
\begin{tabular}{@{} l l @{}}
\toprule
\textbf{Bound} & \textbf{Expression} \\
\midrule
Lower bound & $D$ \\
Upper bound & $D + \dfrac{1}{2} \sum_i n_i P_i$ \\
\bottomrule
\end{tabular}
\end{center}

\begin{tipbox}[Bug 1 vs Bug 2 Comparison]
\begin{center}
\renewcommand{\arraystretch}{1.3}
\begin{tabular}{@{} l p{5.5cm} p{5.5cm} @{}}
\toprule
& \textbf{Bug 1} & \textbf{Bug 2} \\
\midrule
Strategy & Exhaustive (explore all before deciding) & Greedy (take first improvement) \\
Upper bound & $D + \frac{3}{2}\sum P_i$ & $D + \frac{1}{2}\sum n_i P_i$ \\
Performance & More predictable & Often better, but can be worse \\
\bottomrule
\end{tabular}
\end{center}

\textbf{Bug 2 beats Bug 1:} When obstacles have few m-line crossings.

\textbf{Bug 1 beats Bug 2:} When obstacles have many m-line crossings (e.g., spiral shapes).
\end{tipbox}

\separator

% --- Tangent Bug ---
\subsubsection{Tangent Bug Algorithm}

Tangent Bug extends Bug 2 to robots with \textbf{finite-range sensors} (not just contact).

\begin{defbox}[Raw Distance Function]
The \term{raw distance function} $\rho: \mathbb{R}^2 \times S^1 \to \mathbb{R}$ gives the distance to the closest obstacle along a ray from position $x$ at angle $\theta$:
\[
\rho(x, \theta) = \min_{\lambda \in (0,\infty)} d\left(x, x + \lambda(\cos\theta, \sin\theta)^T\right)
\]
such that $x + \lambda(\cos\theta, \sin\theta)^T \in \bigcup_i \mathcal{WO}_i$.

The \term{saturated} version $\rho_R(x,\theta)$ clips at sensor range $R$:
\[
\rho_R(x,\theta) = \begin{cases} \rho(x,\theta) & \text{if } \rho(x,\theta) < R \\ \infty & \text{otherwise} \end{cases}
\]
\end{defbox}

\begin{thmbox}[Tangent Bug Strategy]
\textbf{Motion-to-goal:}
\begin{itemize}
  \item Move toward goal if path is clear.
  \item If obstacles visible, choose discontinuity point $O_i$ that minimizes heuristic: $d(x, O_i) + d(O_i, q_{\text{goal}})$.
  \item Switch to boundary-following if heuristic distance starts increasing.
\end{itemize}

\textbf{Boundary-following:}
\begin{itemize}
  \item Maintain $d_{\text{followed}}$: shortest distance from sensed boundary to goal.
  \item Maintain $d_{\text{reach}}$: shortest distance from blocking obstacle to goal.
  \item Switch back to motion-to-goal when $d_{\text{reach}} < d_{\text{followed}}$.
\end{itemize}
\end{thmbox}

\begin{tipbox}[Tangent Bug: Effect of Sensor Range]
\begin{itemize}
  \item \textbf{Zero range} ($R = 0$): Degenerates to Bug 2 behavior.
  \item \textbf{Finite range}: Can ``look ahead'' and make smarter decisions.
  \item \textbf{Infinite range}: Optimal path planning (full visibility).
\end{itemize}
\end{tipbox}

\bigseparator

% =============================================================================
%                         GEOMETRY REPRESENTATION
% =============================================================================
\subsection{Geometry Representation}

To plan motion, we need to represent:
\begin{itemize}
  \item The \textbf{robot}: set of controllable rigid bodies.
  \item The \textbf{obstacles}: fixed rigid bodies the robot cannot collide with.
\end{itemize}

\separator

\subsubsection{Polygons and Polyhedra}

\begin{defbox}[Polygon / Polyhedron]
\begin{itemize}
  \item \term{Polygon} (2D): described by ordered sequence of vertices or line segments.
  \item \term{Polyhedron} (3D): described by vertices, edges, and faces.
\end{itemize}
Types: triangles, convex polygons, simple polygons, general polygons.
\end{defbox}

\separator

\subsubsection{Convexity}

\begin{defbox}[Convex Set]
A subset $X \subseteq \mathbb{R}^n$ is \term{convex} iff for any two points $x_1, x_2 \in X$, the entire line segment between them lies in $X$:
\[
\lambda x_1 + (1-\lambda) x_2 \in X \quad \forall x_1, x_2 \in X, \; \forall \lambda \in [0,1]
\]
A set that is not convex is called \term{non-convex}.
\end{defbox}

\begin{defbox}[Convex Hull]
The \term{convex hull} of a set of points $\{v_1, \ldots, v_m\}$ is the smallest convex set containing all points:
\[
X = \text{ConvH}(v_1, v_2, \ldots, v_m)
\]
For a polytope in $\mathbb{R}^n$, we need $m \geq n+1$ vertices.
\end{defbox}

\begin{tipbox}[Representing Non-Convex Sets]
Non-convex polygons can be represented as a \textbf{union of convex polygons}:
\[
X = X_1 \cup X_2 \cup \cdots \cup X_n
\]
\textbf{Note:} Decomposing a non-convex set into convex pieces is NP-hard in general!
\end{tipbox}

\separator

\subsubsection{Primitives and Half-Spaces}

\begin{defbox}[Primitive]
A \term{primitive} $H$ is defined by an implicit function $f: \mathbb{R}^n \to \mathbb{R}$:
\[
H = \{x \in \mathbb{R}^n \mid f(x) \leq 0\}
\]
\begin{itemize}
  \item \textbf{Linear primitive} (half-plane/half-space): $f(x) = ax + by + c$ (2D).
  \item \textbf{Algebraic primitive}: $f$ is a polynomial (e.g., $f(x,y) = x^2 + y^2 - r^2$).
\end{itemize}
\end{defbox}

\begin{thmbox}[Convex Polytopes via Half-Spaces]
A convex polygon (polyhedron) can be described as the \textbf{intersection} of half-planes (half-spaces):
\[
X = H_1 \cap H_2 \cap \cdots \cap H_m
\]
For a polytope in $\mathbb{R}^n$, we need $m \geq n+1$ half-spaces.
\end{thmbox}

\begin{exbox}[Worked Example: Hexagon via Half-Planes]
To define a regular hexagon centered at origin, use 6 half-planes:
\[
H_i: a_i x + b_i y + c_i \leq 0, \quad i = 1, \ldots, 6
\]
where each $(a_i, b_i)$ is the outward normal to edge $i$.
\end{exbox}

\separator

\subsubsection{Semi-Algebraic Sets}

\begin{defbox}[Semi-Algebraic Set]
A \term{semi-algebraic set} is constructed using finite unions and intersections of algebraic primitives (sets defined by polynomial inequalities).
\end{defbox}

\begin{exbox}[Example: Smiley Face]
Define a smiley using:
\begin{align*}
H_{\text{head}}&: x^2 + y^2 - r_1^2 \leq 0 \quad \text{(head circle)} \\
H_{\text{Leye}}&: -(x-x_2)^2 - (y-y_2)^2 + r_2^2 \leq 0 \quad \text{(left eye hole)} \\
H_{\text{Reye}}&: -(x-x_3)^2 - (y-y_3)^2 + r_3^2 \leq 0 \quad \text{(right eye hole)} \\
H_{\text{mouth}}&: -\frac{x^2}{a^2} - \frac{(y-y_4)^2}{b^2} + 1 \leq 0 \quad \text{(mouth ellipse)}
\end{align*}
The smiley is: $X = H_{\text{head}} \cap H_{\text{Leye}} \cap H_{\text{Reye}} \cap H_{\text{mouth}}$
\end{exbox}

\separator

\subsubsection{Point Clouds and Octomaps}

\begin{defbox}[Point Cloud]
A \term{point cloud} is a set of 3D points, typically acquired from sensors (LIDAR, depth cameras). Each point represents a surface sample.
\end{defbox}

\begin{defbox}[Voxel]
A \term{voxel} (volumetric pixel) is a 3D grid cell used in occupancy grids.
\end{defbox}

\begin{defbox}[Octomap]
An \term{Octomap} is a probabilistic 3D occupancy map based on \term{octrees}:
\begin{itemize}
  \item Efficiently represents free, occupied, and \textbf{unknown} space.
  \item Supports multi-resolution representation.
  \item Can fuse multiple sensor readings probabilistically.
\end{itemize}
\end{defbox}

\begin{tipbox}[Why Octomaps?]
\begin{itemize}
  \item \textbf{Memory efficient:} Hierarchical structure; empty/uniform regions use little memory.
  \item \textbf{Handles uncertainty:} Probabilistic occupancy values.
  \item \textbf{Updatable:} New sensor data can be fused incrementally.
\end{itemize}
\end{tipbox}

\bigseparator

\subsection{Summary Table: Geometry Representations}

\begin{center}
\renewcommand{\arraystretch}{1.3}
\begin{tabular}{@{} l p{4cm} p{4cm} p{3cm} @{}}
\toprule
\textbf{Representation} & \textbf{Description} & \textbf{Pros} & \textbf{Cons} \\
\midrule
Boundary points & Vertices defining polygon & Simple, exact & Limited to polygons \\
Convex hull & Smallest convex set containing points & Easy intersection tests & Only convex shapes \\
Half-spaces & Intersection of linear inequalities & Exact, good for collision & Only convex polytopes \\
Meshes & Triangulated surfaces & Flexible, standard format & No volume info \\
Algebraic sets & Polynomial inequalities & Smooth curves/surfaces & Complex computations \\
Point clouds & Raw sensor samples & Direct from sensors & Unstructured, noisy \\
Octomaps & Hierarchical voxel grid & Efficient, probabilistic & Discrete resolution \\
\bottomrule
\end{tabular}
\end{center}

\bigseparator

\subsection{Worked Examples}

\begin{exbox}[Example 1: Bug 1 Path Analysis]
\textbf{Problem:} A point robot uses Bug 1 to navigate from $q_{\text{start}}$ to $q_{\text{goal}}$ with straight-line distance $D = 10$ units. There are two obstacles with perimeters $P_1 = 6$ and $P_2 = 8$. What are the path length bounds?

\textbf{Solution:}
\begin{itemize}
  \item Lower bound: $D = 10$ (direct path if no obstacles block).
  \item Upper bound: $D + \frac{3}{2}(P_1 + P_2) = 10 + \frac{3}{2}(6 + 8) = 10 + 21 = 31$ units.
\end{itemize}
\end{exbox}

\begin{exbox}[Example 2: Bug 2 vs Bug 1]
\textbf{Problem:} An obstacle has perimeter $P = 20$ and is crossed by the m-line $n = 8$ times. Compare worst-case path contribution for Bug 1 vs Bug 2.

\textbf{Solution:}
\begin{itemize}
  \item Bug 1 contribution: $\frac{3}{2} \times 20 = 30$ units.
  \item Bug 2 contribution: $\frac{1}{2} \times 8 \times 20 = 80$ units.
\end{itemize}
Bug 1 performs better here because the obstacle has many m-line crossings.
\end{exbox}

\begin{exbox}[Example 3: Convexity Check]
\textbf{Problem:} Is the set $X = \{(x,y) : x^2 + y^2 \leq 1\}$ convex?

\textbf{Solution:} Yes. For any two points $p_1, p_2$ on or inside the unit circle, the line segment between them lies entirely within the circle. This follows from the triangle inequality.
\end{exbox}

\begin{exbox}[Example 4: Triangle via Half-Planes]
\textbf{Problem:} Represent the triangle with vertices $(0,0)$, $(4,0)$, $(2,3)$ using half-planes.

\textbf{Solution:} Find the three edge equations:
\begin{itemize}
  \item Edge 1 ($(0,0)$ to $(4,0)$): $y \geq 0 \Rightarrow -y \leq 0$
  \item Edge 2 ($(4,0)$ to $(2,3)$): $3x + 2y \leq 12$
  \item Edge 3 ($(2,3)$ to $(0,0)$): $3x - 2y \geq 0 \Rightarrow -3x + 2y \leq 0$
\end{itemize}
Triangle: $H_1 \cap H_2 \cap H_3$.
\end{exbox}
